\section{Preparação dos dados}

\begin{frame}[fragile,t,label=cd-preparacao]{3. Qualidade dos dados}
Após conferir os tipos das colunas, verificam-se registros repetidos, ausentes (\texttt{NaN}) e valores fora dos limites esperados.
\EMASCcodelabel{Exemplo prático}
\begin{lstlisting}
print("Duplicatas:", df.duplicated().sum())
print("IDs repetidos:", df["id"].duplicated().sum())
print("Ausentes:\n", df.isna().sum())
fora = df["nota"].notna() & ~df["nota"].between(0, 10)
print("Notas fora de 0 a 10:\n", df[fora])
\end{lstlisting}
\begin{EMASCnote}
\texttt{sum()} conta os valores \texttt{True} como 1. Há duas notas ausentes; a tabela vazia indica que nenhuma nota registrada está fora de 0 a 10.
\end{EMASCnote}
\end{frame}

\begin{frame}[fragile,t]{3.1. Tratamento de valores ausentes}
Existem diferentes formas de tratar valores ausentes. A escolha depende do significado desses valores e do objetivo da análise.
\EMASCcodelabel{Exemplo prático}
\begin{lstlisting}
pendentes = df[df["nota"].isna()]
print(pendentes[["id", "curso"]])

df_com_notas = df.dropna(subset=["nota"]).copy()
print("Linhas e colunas:", df_com_notas.shape)
\end{lstlisting}
\begin{EMASCnote}
Uma nota ausente não significa nota zero. Os 18 registros ficam em \texttt{df}; \texttt{dropna()} mantém apenas os 16 com nota na cópia.
\end{EMASCnote}
\small
\texttt{fillna(valor)} é outra possibilidade, quando a substituição puder ser justificada. Neste problema, substituir notas inventaria informações.
\end{frame}

\begin{frame}[fragile,t]{3.2. Criação de colunas}
Novas colunas podem ser calculadas ou obtidas por meio de uma função. Vamos converter a frequência em porcentagem e classificar as notas.
\EMASCcodelabel{Exemplo prático}
\begin{lstlisting}
df_com_notas["frequencia_pct"] = df_com_notas["frequencia"] * 100
def classificar(nota):
    if nota >= 7:
        return "Nota >= 7"
    return "Nota < 7"
df_com_notas["situacao"] = df_com_notas["nota"].apply(classificar)
print(df_com_notas[["frequencia_pct", "nota", "situacao"]].head())
\end{lstlisting}
\begin{EMASCnote}
\texttt{apply()} aplica a função a cada nota. O limite 7 é didático, não uma regra oficial de aprovação.
\end{EMASCnote}
\end{frame}
